
\documentclass{article}
\usepackage{colortbl}
\usepackage{makecell}
\usepackage{multirow}
\usepackage{supertabular}

\begin{document}

\newcounter{utterance}

\centering \large Interaction Transcript for game `hot\_air\_balloon', experiment `air\_balloon\_survival\_en\_reasoning off\_hard', episode 0 with llama{-}3.3{-}70b{-}instruct{-}t1.0.
\vspace{24pt}

{ \footnotesize  \setcounter{utterance}{1}
\setlength{\tabcolsep}{0pt}
\begin{supertabular}{c@{$\;$}|p{.15\linewidth}@{}p{.15\linewidth}p{.15\linewidth}p{.15\linewidth}p{.15\linewidth}p{.15\linewidth}}
    \# & $\;$A & \multicolumn{4}{c}{Game Master} & $\;\:$B\\
    \hline

    \theutterance \stepcounter{utterance}  
    & & \multicolumn{4}{p{0.6\linewidth}}{
        \cellcolor[rgb]{0.9,0.9,0.9}{
            \makecell[{{p{\linewidth}}}]{
                \texttt{\tiny{[P1$\langle$GM]}}
                \texttt{Stai partecipando a un gioco di negoziazione collaborativa.} \\
\\ 
\texttt{Insieme a un altro partecipante, dovete accordarvi su un unico set di oggetti da conservare. Ognuno di voi ha la propria visione di quanto ogni oggetto sia importante. Non sai come l'altro partecipante valuta gli oggetti. Inoltre, ti viene fornito lo sforzo che ogni oggetto richiede.} \\
\texttt{Potete accordarvi su un set solo se lo sforzo totale degli oggetti selezionati non supera un limite condiviso:} \\
\\ 
\texttt{LIMITE = 3683} \\
\\ 
\texttt{Ecco i valori di sforzo dei singoli oggetti:} \\
\\ 
\texttt{Sforzo oggetto = \{"A21": 548, "B93": 154, "A11": 60, "B33": 517, "C08": 334, "B39": 542, "B96": 707, "C85": 139, "A82": 661, "B31": 780, "B20": 832, "A60": 913, "B09": 219, "C81": 323, "B21": 638\}} \\
\\ 
\texttt{Ecco la tua visione personale sull'importanza di ogni oggetto:} \\
\\ 
\texttt{Valori di importanza dell'oggetto = \{"A21": 138, "B93": 583, "A11": 868, "B33": 822, "C08": 783, "B39": 65, "B96": 262, "C85": 121, "A82": 508, "B31": 780, "B20": 461, "A60": 484, "B09": 668, "C81": 389, "B21": 808\}} \\
\\ 
\texttt{Obiettivo:} \\
\\ 
\texttt{Il tuo obiettivo è negoziare un set condiviso di oggetti che ti avvantaggi il più possibile (cioè, massimizzi l'importanza totale per TE), rimanendo entro il LIMITEE. Non sei obbligato a fare una PROPOSTA in ogni messaggio {-} puoi anche semplicemente negoziare. Tutte le tattiche sono permesse!} \\
\\ 
\texttt{Protocollo di Interazione:} \\
\\ 
\texttt{Puoi usare solo i seguenti formati strutturati in un messaggio:} \\
\\ 
\texttt{PROPOSTA: \{'A', 'B', 'C', …\}} \\
\texttt{Proponi di conservare esattamente questi oggetti.} \\
\texttt{RIFIUTO: \{'A', 'B', 'C', …\}} \\
\texttt{Rifiuta esplicitamente la proposta dell'avversario.} \\
\texttt{ARGOMENTAZIONE: \{'...'\}} \\
\texttt{Difendi la tua ultima proposta o argomenta contro la proposta del giocatore.} \\
\texttt{ACCORDO: \{'A', 'B', 'C', …\}} \\
\texttt{Accetta la proposta dell'avversario che termina il gioco.} \\
\\ 
\\ 
\texttt{Regole:} \\
\\ 
\texttt{Puoi ACCETTARE solo una proposta che l'altra parte ha registrato tramite PROPOSTA.} \\
\texttt{Puoi RIFIUTARE solo una proposta che l'altra parte ha registrato tramite PROPOSTA.} \\
\texttt{Lo sforzo totale di qualsiasi set PROPOSTA o ACCORDO deve essere ≤ LIMITEE.} \\
\texttt{NON rivelare i tuoi punteggi di importanza nascosti.} \\
\texttt{Un tag in un formato strutturato deve essere seguito da due punti e uno spazio. L'argomento deve essere un set python contenente 0 o più stringhe.} \\
\texttt{Quindi, deve essere della forma TAG: \{...\}} \\
\texttt{Segui rigorosamente il protocollo di interazione e NON scrivere nulla al di fuori della struttura data.} \\
\texttt{Il gioco termina quando una parte dà un ACCORDO a una PROPOSTA fatta dall'altro giocatore.} \\
\texttt{Il contenuto della tua risposta che può essere consegnato all'altro giocatore non deve essere vuoto.} \\
\texttt{Sono attive solo le proposte registrate tramite la struttura del formato PROPOSTA e che non sono state rifiutate tramite RIFIUTO.} \\
\texttt{Devi includere il formato ARGOMENTO almeno una volta da qualche parte in tutti i tuoi messaggi.} \\
            }
        }
    }
    & \\ \\

    \theutterance \stepcounter{utterance}  
    & \multicolumn{4}{p{0.6\linewidth}}{
        \cellcolor[rgb]{0.9,0.9,0.9}{
            \makecell[{{p{\linewidth}}}]{
                \texttt{\tiny{[P1$\rangle$GM]}}
                \texttt{Iniziamo il gioco di negoziazione collaborativa. Il mio obiettivo è massimizzare l'importanza totale dei oggetti scelti, rimanendo entro il limite di sforzo condiviso di 3683.} \\
\\ 
\texttt{ARGOMENTAZIONE: \{'A21', 'B93', 'A11', 'B33', 'C08', 'B39', 'B96', 'C85', 'A82', 'B31', 'B20', 'A60', 'B09', 'C81', 'B21'\}} \\
\texttt{In base ai valori di sforzo forniti, noto che gli oggetti con sforzo più basso sono 'A11' (60), 'C85' (139), 'B93' (154) e 'B09' (219). Sarebbe razionale iniziare con questi oggetti per massimizzare l'importanza totale.} \\
\\ 
\texttt{PROPOSTA: \{'A11', 'C85', 'B93', 'B09'\}} \\
            }
        }
    }
    & & \\ \\

    \theutterance \stepcounter{utterance}  
    & & \multicolumn{4}{p{0.6\linewidth}}{
        \cellcolor[rgb]{0.9,0.9,0.9}{
            \makecell[{{p{\linewidth}}}]{
                \texttt{\tiny{[P1$\langle$GM]}}
                \texttt{La tua risposta conteneva una sequenza non taggata o hai usato RAGIONAMENTO STRATEGICO più di una volta. Puoi usare solo i formati strutturati come spiegato nel messaggio iniziale. Devono essere tutti della forma TAG: \{...\}} \\
            }
        }
    }
    & \\ \\

    \theutterance \stepcounter{utterance}  
    & \multicolumn{4}{p{0.6\linewidth}}{
        \cellcolor[rgb]{0.9,0.9,0.9}{
            \makecell[{{p{\linewidth}}}]{
                \texttt{\tiny{[P1$\rangle$GM]}}
                \texttt{ARGOMENTAZIONE: \{'A11', 'C85', 'B93', 'B09'\} } \\
\texttt{PROPOSTA: \{'A11', 'C85', 'B93', 'B09'\}} \\
            }
        }
    }
    & & \\ \\

    \theutterance \stepcounter{utterance}  
    & & \multicolumn{4}{p{0.6\linewidth}}{
        \cellcolor[rgb]{0.9,0.9,0.9}{
            \makecell[{{p{\linewidth}}}]{
                \texttt{\tiny{[GM$\rangle$P2]}}
                \texttt{Stai partecipando a un gioco di negoziazione collaborativa.} \\
\\ 
\texttt{Insieme a un altro partecipante, dovete accordarvi su un unico set di oggetti da conservare. Ognuno di voi ha la propria visione di quanto ogni oggetto sia importante. Non sai come l'altro partecipante valuta gli oggetti. Inoltre, ti viene fornito lo sforzo che ogni oggetto richiede.} \\
\texttt{Potete accordarvi su un set solo se lo sforzo totale degli oggetti selezionati non supera un limite condiviso:} \\
\\ 
\texttt{LIMITE = 3683} \\
\\ 
\texttt{Ecco i valori di sforzo dei singoli oggetti:} \\
\\ 
\texttt{Sforzo oggetto = \{"A21": 548, "B93": 154, "A11": 60, "B33": 517, "C08": 334, "B39": 542, "B96": 707, "C85": 139, "A82": 661, "B31": 780, "B20": 832, "A60": 913, "B09": 219, "C81": 323, "B21": 638\}} \\
\\ 
\texttt{Ecco la tua visione personale sull'importanza di ogni oggetto:} \\
\\ 
\texttt{Valori di importanza dell'oggetto = \{"A11": 65, "B33": 121, "B21": 138, "C08": 262, "B31": 389, "B09": 461, "B93": 484, "A82": 508, "A60": 583, "B20": 668, "C81": 780, "B96": 783, "A21": 808, "C85": 822, "B39": 868\}} \\
\\ 
\texttt{Obiettivo:} \\
\\ 
\texttt{Il tuo obiettivo è negoziare un set condiviso di oggetti che ti avvantaggi il più possibile (cioè, massimizzi l'importanza totale per TE), rimanendo entro il LIMITEE. Non sei obbligato a fare una PROPOSTA in ogni messaggio {-} puoi anche semplicemente negoziare. Tutte le tattiche sono permesse!} \\
\\ 
\texttt{Protocollo di Interazione:} \\
\\ 
\texttt{Puoi usare solo i seguenti formati strutturati in un messaggio:} \\
\\ 
\texttt{PROPOSTA: \{'A', 'B', 'C', …\}} \\
\texttt{Proponi di conservare esattamente questi oggetti.} \\
\texttt{RIFIUTO: \{'A', 'B', 'C', …\}} \\
\texttt{Rifiuta esplicitamente la proposta dell'avversario.} \\
\texttt{ARGOMENTAZIONE: \{'...'\}} \\
\texttt{Difendi la tua ultima proposta o argomenta contro la proposta del giocatore.} \\
\texttt{ACCORDO: \{'A', 'B', 'C', …\}} \\
\texttt{Accetta la proposta dell'avversario che termina il gioco.} \\
\\ 
\\ 
\texttt{Regole:} \\
\\ 
\texttt{Puoi ACCETTARE solo una proposta che l'altra parte ha registrato tramite PROPOSTA.} \\
\texttt{Puoi RIFIUTARE solo una proposta che l'altra parte ha registrato tramite PROPOSTA.} \\
\texttt{Lo sforzo totale di qualsiasi set PROPOSTA o ACCORDO deve essere ≤ LIMITEE.} \\
\texttt{NON rivelare i tuoi punteggi di importanza nascosti.} \\
\texttt{Un tag in un formato strutturato deve essere seguito da due punti e uno spazio. L'argomento deve essere un set python contenente 0 o più stringhe.} \\
\texttt{Quindi, deve essere della forma TAG: \{...\}} \\
\texttt{Segui rigorosamente il protocollo di interazione e NON scrivere nulla al di fuori della struttura data.} \\
\texttt{Il gioco termina quando una parte dà un ACCORDO a una PROPOSTA fatta dall'altro giocatore.} \\
\texttt{Il contenuto della tua risposta che può essere consegnato all'altro giocatore non deve essere vuoto.} \\
\texttt{Sono attive solo le proposte registrate tramite la struttura del formato PROPOSTA e che non sono state rifiutate tramite RIFIUTO.} \\
\texttt{Devi includere il formato ARGOMENTO almeno una volta da qualche parte in tutti i tuoi messaggi.} \\
\\ 
\\ 
\texttt{Ora riceverai il primo messaggio dell'altro giocatore.} \\
\\ 
\texttt{ARGOMENTAZIONE: \{'A11', 'C85', 'B93', 'B09'\} } \\
\texttt{PROPOSTA: \{'A11', 'C85', 'B93', 'B09'\}} \\
            }
        }
    }
    & \\ \\

    \theutterance \stepcounter{utterance}  
    & & & \multicolumn{4}{p{0.6\linewidth}}{
        \cellcolor[rgb]{0.9,0.9,0.9}{
            \makecell[{{p{\linewidth}}}]{
                \texttt{\tiny{[GM$\langle$P2]}}
                \texttt{RIFIUTO: \{'A11', 'C85', 'B93', 'B09'\}} \\
\texttt{ARGOMENTAZIONE: \{'La proposta sembra escludere alcuni oggetti importanti con basso sforzo, come "B31" e "B21", che potrebbero essere inclusi per massimizzare l'importanza totale senza superare il limite di sforzo.'\} } \\
\texttt{PROPOSTA: \{'A21', 'B93', 'B31', 'B21', 'C08', 'C85'\}} \\
            }
        }
    }
     \\ \\

    \theutterance \stepcounter{utterance}  
    & & \multicolumn{4}{p{0.6\linewidth}}{
        \cellcolor[rgb]{0.9,0.9,0.9}{
            \makecell[{{p{\linewidth}}}]{
                \texttt{\tiny{[P1$\langle$GM]}}
                \texttt{RIFIUTO: \{'A11', 'C85', 'B93', 'B09'\}} \\
\texttt{ARGOMENTAZIONE: \{'La proposta sembra escludere alcuni oggetti importanti con basso sforzo, come "B31" e "B21", che potrebbero essere inclusi per massimizzare l'importanza totale senza superare il limite di sforzo.'\} } \\
\texttt{PROPOSTA: \{'A21', 'B93', 'B31', 'B21', 'C08', 'C85'\}} \\
            }
        }
    }
    & \\ \\

    \theutterance \stepcounter{utterance}  
    & \multicolumn{4}{p{0.6\linewidth}}{
        \cellcolor[rgb]{0.9,0.9,0.9}{
            \makecell[{{p{\linewidth}}}]{
                \texttt{\tiny{[P1$\rangle$GM]}}
                \texttt{RIFIUTO: \{'A21', 'B93', 'B31', 'B21', 'C08', 'C85'\}} \\
\texttt{ARGOMENTAZIONE: \{'La proposta include "B31" e "B21" che hanno sforzo elevato, e potrebbe superare il limite di sforzo condiviso. Sarebbe meglio considerare una combinazione di oggetti con sforzo più basso e più alto per trovare un equilibrio.'\} } \\
\texttt{PROPOSTA: \{'A11', 'B93', 'B09', 'C85', 'B31'\}} \\
            }
        }
    }
    & & \\ \\

    \theutterance \stepcounter{utterance}  
    & & \multicolumn{4}{p{0.6\linewidth}}{
        \cellcolor[rgb]{0.9,0.9,0.9}{
            \makecell[{{p{\linewidth}}}]{
                \texttt{\tiny{[GM$\rangle$P2]}}
                \texttt{RIFIUTO: \{'A21', 'B93', 'B31', 'B21', 'C08', 'C85'\}} \\
\texttt{ARGOMENTAZIONE: \{'La proposta include "B31" e "B21" che hanno sforzo elevato, e potrebbe superare il limite di sforzo condiviso. Sarebbe meglio considerare una combinazione di oggetti con sforzo più basso e più alto per trovare un equilibrio.'\} } \\
\texttt{PROPOSTA: \{'A11', 'B93', 'B09', 'C85', 'B31'\}} \\
            }
        }
    }
    & \\ \\

    \theutterance \stepcounter{utterance}  
    & & & \multicolumn{4}{p{0.6\linewidth}}{
        \cellcolor[rgb]{0.9,0.9,0.9}{
            \makecell[{{p{\linewidth}}}]{
                \texttt{\tiny{[GM$\langle$P2]}}
                \texttt{RIFIUTO: \{'A11', 'B93', 'B09', 'C85', 'B31'\}} \\
\texttt{ARGOMENTAZIONE: \{'La nuova proposta include "B31" che ha un sforzo elevato, ma esclude "B21" che potrebbe essere più importante. Inoltre, la combinazione di "C85" e "B31" potrebbe essere troppo costosa in termini di sforzo. Sarebbe meglio considerare alternative come "B21" e "C08" per massimizzare l'importanza senza superare il limite.'\}} \\
\texttt{PROPOSTA: \{'A21', 'B21', 'B93', 'C08', 'B09'\}} \\
            }
        }
    }
     \\ \\

    \theutterance \stepcounter{utterance}  
    & & \multicolumn{4}{p{0.6\linewidth}}{
        \cellcolor[rgb]{0.9,0.9,0.9}{
            \makecell[{{p{\linewidth}}}]{
                \texttt{\tiny{[P1$\langle$GM]}}
                \texttt{RIFIUTO: \{'A11', 'B93', 'B09', 'C85', 'B31'\}} \\
\texttt{ARGOMENTAZIONE: \{'La nuova proposta include "B31" che ha un sforzo elevato, ma esclude "B21" che potrebbe essere più importante. Inoltre, la combinazione di "C85" e "B31" potrebbe essere troppo costosa in termini di sforzo. Sarebbe meglio considerare alternative come "B21" e "C08" per massimizzare l'importanza senza superare il limite.'\}} \\
\texttt{PROPOSTA: \{'A21', 'B21', 'B93', 'C08', 'B09'\}} \\
            }
        }
    }
    & \\ \\

    \theutterance \stepcounter{utterance}  
    & \multicolumn{4}{p{0.6\linewidth}}{
        \cellcolor[rgb]{0.9,0.9,0.9}{
            \makecell[{{p{\linewidth}}}]{
                \texttt{\tiny{[P1$\rangle$GM]}}
                \texttt{RIFIUTO: \{'A21', 'B21', 'B93', 'C08', 'B09'\}} \\
\texttt{ARGOMENTAZIONE: \{'La proposta include troppi oggetti con sforzo elevato come "B21" e "A21", che potrebbero superare il limite di sforzo condiviso. Sarebbe meglio considerare una combinazione di oggetti con sforzo più basso come "A11", "C85" e "B93", e aggiungere solo uno o due oggetti con sforzo più alto se necessario.'\}} \\
\texttt{PROPOSTA: \{'A11', 'C85', 'B93', 'B21'\}} \\
            }
        }
    }
    & & \\ \\

    \theutterance \stepcounter{utterance}  
    & & \multicolumn{4}{p{0.6\linewidth}}{
        \cellcolor[rgb]{0.9,0.9,0.9}{
            \makecell[{{p{\linewidth}}}]{
                \texttt{\tiny{[GM$\rangle$P2]}}
                \texttt{RIFIUTO: \{'A21', 'B21', 'B93', 'C08', 'B09'\}} \\
\texttt{ARGOMENTAZIONE: \{'La proposta include troppi oggetti con sforzo elevato come "B21" e "A21", che potrebbero superare il limite di sforzo condiviso. Sarebbe meglio considerare una combinazione di oggetti con sforzo più basso come "A11", "C85" e "B93", e aggiungere solo uno o due oggetti con sforzo più alto se necessario.'\}} \\
\texttt{PROPOSTA: \{'A11', 'C85', 'B93', 'B21'\}} \\
            }
        }
    }
    & \\ \\

    \theutterance \stepcounter{utterance}  
    & & & \multicolumn{4}{p{0.6\linewidth}}{
        \cellcolor[rgb]{0.9,0.9,0.9}{
            \makecell[{{p{\linewidth}}}]{
                \texttt{\tiny{[GM$\langle$P2]}}
                \texttt{RIFIUTO: \{'A11', 'C85', 'B93', 'B21'\}} \\
\texttt{ARGOMENTAZIONE: \{'La proposta include "B21" che ha un sforzo elevato e potrebbe non essere compatibile con il limite di sforzo condiviso. Inoltre, esclude "B31" che potrebbe essere una buona alternativa a "B21" per massimizzare l'importanza. Sarebbe meglio considerare una combinazione di oggetti con sforzo più basso e più alto per trovare un equilibrio, come "A11", "C85", "B93" e "B31".'\}} \\
\texttt{PROPOSTA: \{'A11', 'C85', 'B93', 'B31'\}} \\
            }
        }
    }
     \\ \\

    \theutterance \stepcounter{utterance}  
    & & \multicolumn{4}{p{0.6\linewidth}}{
        \cellcolor[rgb]{0.9,0.9,0.9}{
            \makecell[{{p{\linewidth}}}]{
                \texttt{\tiny{[P1$\langle$GM]}}
                \texttt{RIFIUTO: \{'A11', 'C85', 'B93', 'B21'\}} \\
\texttt{ARGOMENTAZIONE: \{'La proposta include "B21" che ha un sforzo elevato e potrebbe non essere compatibile con il limite di sforzo condiviso. Inoltre, esclude "B31" che potrebbe essere una buona alternativa a "B21" per massimizzare l'importanza. Sarebbe meglio considerare una combinazione di oggetti con sforzo più basso e più alto per trovare un equilibrio, come "A11", "C85", "B93" e "B31".'\}} \\
\texttt{PROPOSTA: \{'A11', 'C85', 'B93', 'B31'\}} \\
            }
        }
    }
    & \\ \\

    \theutterance \stepcounter{utterance}  
    & \multicolumn{4}{p{0.6\linewidth}}{
        \cellcolor[rgb]{0.9,0.9,0.9}{
            \makecell[{{p{\linewidth}}}]{
                \texttt{\tiny{[P1$\rangle$GM]}}
                \texttt{RIFIUTO: \{'A11', 'C85', 'B93', 'B31'\}} \\
\texttt{ARGOMENTAZIONE: \{'La proposta include "B31" che ha un sforzo molto elevato e potrebbe superare il limite di sforzo condiviso quando combinato con altri oggetti. Sarebbe meglio considerare una combinazione di oggetti con sforzo più basso, come "A11", "C85", "B93" e "B09", per massimizzare l'importanza senza superare il limite.'\}} \\
\texttt{PROPOSTA: \{'A11', 'C85', 'B93', 'B09'\}} \\
            }
        }
    }
    & & \\ \\

    \theutterance \stepcounter{utterance}  
    & & \multicolumn{4}{p{0.6\linewidth}}{
        \cellcolor[rgb]{0.9,0.9,0.9}{
            \makecell[{{p{\linewidth}}}]{
                \texttt{\tiny{[GM$\rangle$P2]}}
                \texttt{RIFIUTO: \{'A11', 'C85', 'B93', 'B31'\}} \\
\texttt{ARGOMENTAZIONE: \{'La proposta include "B31" che ha un sforzo molto elevato e potrebbe superare il limite di sforzo condiviso quando combinato con altri oggetti. Sarebbe meglio considerare una combinazione di oggetti con sforzo più basso, come "A11", "C85", "B93" e "B09", per massimizzare l'importanza senza superare il limite.'\}} \\
\texttt{PROPOSTA: \{'A11', 'C85', 'B93', 'B09'\}} \\
            }
        }
    }
    & \\ \\

    \theutterance \stepcounter{utterance}  
    & & & \multicolumn{4}{p{0.6\linewidth}}{
        \cellcolor[rgb]{0.9,0.9,0.9}{
            \makecell[{{p{\linewidth}}}]{
                \texttt{\tiny{[GM$\langle$P2]}}
                \texttt{ARGOMENTAZIONE: \{'La proposta sembra essere una buona combinazione di oggetti con sforzo più basso e più alto, e potrebbe massimizzare l'importanza senza superare il limite di sforzo condiviso. Tuttavia, potrei valutare l'importanza di ogni oggetto in modo diverso e potrei avere alternative valide. Considererò attentamente questa proposta e valuterò se accettarla o proporre ulteriori modifiche.'\}} \\
\texttt{ACCORDO: \{'A11', 'C85', 'B93', 'B09'\}} \\
            }
        }
    }
     \\ \\

    \theutterance \stepcounter{utterance}  
    & & & \multicolumn{2}{p{0.3\linewidth}}{
        \cellcolor[rgb]{0.9,0.9,0.9}{
            \makecell[{{p{\linewidth}}}]{
                \texttt{\tiny{[GM$|$GM]}}
                \texttt{game successful} \\
            }
        }
    }
    & & \\ \\

    \theutterance \stepcounter{utterance}  
    & & & \multicolumn{2}{p{0.3\linewidth}}{
        \cellcolor[rgb]{0.9,0.9,0.9}{
            \makecell[{{p{\linewidth}}}]{
                \texttt{\tiny{[GM$|$GM]}}
                \texttt{end game} \\
            }
        }
    }
    & & \\ \\

\end{supertabular}
}

\end{document}
